% https://tex.stackexchange.com/a/729888

\documentclass[tikz,border=8pt]{standalone}
\usetikzlibrary{positioning}
\usetikzlibrary{chains}
\makeatletter
% copy from: https://github.com/pgf-tikz/pgf/blob/master/tex/generic/pgf/frontendlayer/tikz/libraries/tikzlibrarypositioning.code.tex – simpler than trying to \let them to another key
\tikzset{orig above  left/.code=\tikz@lib@place@handle@{#1}{south east}{-1}{ 1}{north west}{0.707106781}}%
\tikzset{orig above right/.code=\tikz@lib@place@handle@{#1}{south west}{ 1}{ 1}{north east}{0.707106781}}%
\tikzset{orig below  left/.code=\tikz@lib@place@handle@{#1}{north east}{-1}{-1}{south west}{0.707106781}}%
\tikzset{orig below right/.code=\tikz@lib@place@handle@{#1}{north west}{ 1}{-1}{south east}{0.707106781}}%
% I'm using the low-level PGFkeys macros since I'm using \pgfeov internally – with a few auxilliary macros this could easily be done in /key/.format=…
\pgfkeysdef{/tikz/above right}{\pgfutil@ifnextchar\bgroup{\pgfkeysvalueof{/explorer/above right/.@cmd}}{\pgfkeysvalueof{/tikz/orig above right/.@cmd}}#1\pgfeov}
\pgfkeysdef{/tikz/above  left}{\pgfutil@ifnextchar\bgroup{\pgfkeysvalueof{/explorer/above  left/.@cmd}}{\pgfkeysvalueof{/tikz/orig above  left/.@cmd}}#1\pgfeov}
\pgfkeysdef{/tikz/below  left}{\pgfutil@ifnextchar\bgroup{\pgfkeysvalueof{/explorer/below  left/.@cmd}}{\pgfkeysvalueof{/tikz/orig below  left/.@cmd}}#1\pgfeov}
\pgfkeysdef{/tikz/below right}{\pgfutil@ifnextchar\bgroup{\pgfkeysvalueof{/explorer/below right/.@cmd}}{\pgfkeysvalueof{/tikz/orig below right/.@cmd}}#1\pgfeov}
\pgfkeysdefnargs{/explorer/above right}{2}{\pgfmathsincos{#1}\tikz@lib@place@handle@{#2}{180+#1}{\pgfmathresultx}{\pgfmathresulty}{#1}{1}}
\pgfkeysdefnargs{/explorer/above  left}{2}{\pgfmathsincos{180-(#1)}\tikz@lib@place@handle@{#2}{#1}{\pgfmathresultx}{\pgfmathresulty}{180-(#1)}{1}}
\pgfkeysdefnargs{/explorer/below  left}{2}{\pgfmathsincos{180+#1}\tikz@lib@place@handle@{#2}{#1}{\pgfmathresultx}{\pgfmathresulty}{180+#1}{1}}
\pgfkeysdefnargs{/explorer/below right}{2}{\pgfmathsincos{-(#1)}\tikz@lib@place@handle@{#2}{180-(#1)}{\pgfmathresultx}{\pgfmathresulty}{-(#1)}{1}}
\pgfkeysdefnargs{/tikz/angle}{2}{\pgfmathsincos{#1}\tikz@lib@place@handle@{#2}{180+#1}{\pgfmathresultx}{\pgfmathresulty}{#1}{1}}
\makeatother

\begin{document}
\begin{tikzpicture}[every node/.style={draw, circle, inner sep=5pt, outer sep=0pt, fill=cyan!20}]
\node (O) {$1$};
\node (A) [below  left={60}{of O}] {$2$} edge (O);
\node (B) [below right={30}{of A}] {$3$} edge (A);
\end{tikzpicture}

\begin{tikzpicture}[
  every node/.style={draw, circle, inner sep=5pt, outer sep=0pt, fill=cyan!20},
  start chain=going {below \ifodd\tikzchaincount right={60}\else left={30}\fi{of \tikzchainprevious}},
  every on chain/.style={join}]
\foreach \v in {1, ..., 9} \node[on chain] {\v};
\end{tikzpicture}

\begin{tikzpicture}[
  every node/.style={draw, circle, inner sep=5pt, outer sep=0pt, fill=cyan!20, minimum size=1cm},
  bl/.style={below left={60}{of #1}},
  br/.style={below right={30}{of #1}}
]
\node       (1) {1};
\node[bl=1] (2) { } node[bl=2] (3) {3};
\node[br=1] (4) { } node[br=4] (5) {5};
\node[left=of 5] (1') {1};
\draw (1) -- (2) -- (3) -- (1') -- (5) -- (4) -- (1);
\draw[red] (1) -- (1');
\end{tikzpicture}

\tikz[node distance=1.5cm and 3cm, nodes={draw, circle}, on grid]
 \draw ellipse [x radius=3cm, y radius=1.5cm] node (O) {O}
   node foreach \ang in {0, 30, ..., 359}[angle={\ang}{of O}]{\ang};
\end{document}
